
\section{Each type-III subtype has the commuting language}\label{sec:typeIII}
The state and host may both be chosen. This quantifier is essential.

\begin{backgroundlemma}[Expected-inclusion transfer; \cite{Takesaki,ReadType}]
If $E:N\to M$ is a faithful normal conditional expectation, every
faithful standard-form correlation on $M$ is realized on $N$.
\end{backgroundlemma}
\begin{proof}
For the state $\varphi_N=\varphi_M\circ E$, put
$\mathcal H_{\rm inc}=\overline{M\xi_{\varphi_N}}\subset L^2(N)$.
The map $x\xi_{\varphi_M}\mapsto x\xi_{\varphi_N}$ is an isometry.
Takesaki's expectation theorem \cite{Takesaki} makes $M$ invariant
under the modular group. The Tomita operators agree on the common
core $M\xi$, and their conjugations restrict as $\mathcal J_N|_{\mathcal H_{\rm inc}}=\mathcal J_M$.
Thus for left effects $A$ and right effects $\mathcal J_MB\mathcal J_M$,
\[
 \langle\xi_{\varphi_N},A \mathcal J_NB\mathcal J_N\xi_{\varphi_N}\rangle
 =\langle\xi_{\varphi_M},A \mathcal J_MB\mathcal J_M\xi_{\varphi_M}\rangle.
\]
Positivity, sums and all probabilities are preserved. Support
reduction handles nonfaithful states. This is also
\cite[Lemma 3.1]{ReadType}.
\end{proof}

\begin{lemma}\label{lem:absorb}
For any factor $N_{\rm in}$ with separable predual, $N_{\rm in}\vN R_\infty$ is
$\III_1$, and contains $N_{\rm in}$ with a faithful normal expectation.
\end{lemma}
\begin{proof}
The expectation is $\operatorname{id}\otimes\rho$ for a faithful normal
state $\rho$ on $R_\infty$.
We use the proof, rather than the stronger hyperfinite conclusion,
of \cite[Lemma 3.8]{HM}.
For each $0<\lambda<1$, choose $0<\mu<1$ with
$\log\lambda/\log\mu$ irrational.
The injective factors satisfy
$R_\infty\cong R_\lambda\vN R_\mu$ and
$R_\infty\vN R_\lambda\cong R_\infty$.
Consequently $(N_{\rm in}\vN R_\infty)\vN R_\lambda
\cong N_{\rm in}\vN R_\infty$.
The asymptotic-ratio argument of Connes \cite[Section 3.6]{Connes}, used in that proof, gives
$\lambda\in r_\infty(N_{\rm in}\vN R_\infty)$.
Connes's implication \cite[Theorem 3.6.1]{Connes} gives
$r_\infty\setminus\{0\}\subseteq S$.
Every $0<\lambda<1$ therefore belongs to the $S$-invariant, and the
factor is $\III_1$. These steps do not assume hyperfiniteness of $N_{\rm in}$.
Only the last identification with the unique injective $\III_1$
factor in Haagerup--Musat needs hyperfiniteness, and we do not make
that identification here.
\end{proof}

\begin{theorem}\label{thm:III}
For each fixed $\lambda\in[0,1]$,
$\perf(\III_\lambda)=\perf_{\co}$.
\end{theorem}
\begin{proof}
The inclusion into commuting perfection is immediate.
For the converse, standardize a commuting-perfect strategy and
choose a factor host $N_{\rm in}$ by central decomposition as in
Section~\ref{sec:intro}. Put $N_{\rm abs}=N_{\rm in}\vN R_\infty$.
By expected-inclusion transfer and Lemma~\ref{lem:absorb}, this $\III_1$
factor hosts the game. Retain it for $\lambda=1$.

Fix a faithful normal state $\varphi$ on $N_{\rm abs}$.
For a countable subgroup $\Gamma\subset\mathbb R$, given the discrete
topology, form
$N_\Gamma=N_{\rm abs}\rtimes_{\sigma^\varphi|_\Gamma}\Gamma$.
The Fourier expectation onto $N_{\rm abs}$ is faithful normal, so transfers
the game. It is a factor: $\mathrm T(N_{\rm abs})=\{0\}$ for a $\III_1$ factor, hence
every nonzero modular time is outer.
If a central element has Fourier coefficient $x_g$, commutation
with $N_{\rm abs}$ makes $x_g$ an intertwiner between $\operatorname{id}$ and
$\sigma_g^\varphi$. A nonzero such coefficient for $g\ne0$ would,
by polar decomposition in a factor, implement that automorphism
innerly. Thus only $g=0$ remains, and $x_0$ is scalar.
Faithfulness of Fourier coefficients proves factoriality.

The continuous core is the crossed product by the modular
$\mathbb R$-action; its dual action restricted to the center is the
\emph{flow of weights}. For factors this flow distinguishes the
type-III subtypes \cite{ConnesTakesaki}.
Here is the core computation, with the dual flow specified.
The modular group of the dual state fixes the canonical unitaries
$u_g$ of the discrete crossed product. Writing $\lambda(t)$ for
the continuous-core implementing unitaries, it follows that
\begin{equation}\label{eq:core}
 \begin{split}
 c(N_\Gamma)
 &\cong c(N_{\rm abs})\rtimes_{\Ad\lambda(g)}\Gamma\\
 &\cong c(N_{\rm abs})\vN L(\Gamma).
 \end{split}
\end{equation}
For the second isomorphism the unitaries
$w_g=\lambda(g)^*u_g$ commute with $c(N_{\rm abs})$ and satisfy the group
law of $\Gamma$; the inverse sends $u_g$ to $\lambda(g)\otimes w_g$.
Since $N_{\rm abs}$ is $\III_1$, $c(N_{\rm abs})$ is a factor, so the center is
$L(\Gamma)=L^\infty(\widehat\Gamma)$.
The dual action obeys
$\theta_s(\lambda(t))=e^{-ist}\lambda(t)$ and
$\theta_s(u_g)=u_g$; therefore
$\theta_s(w_g)=e^{isg}w_g$.
Its action on the center is translation by the character
$\chi_s(g)=e^{isg}$.

For $0<\lambda<1$, take
$\Gamma=t_\lambda\mathbb Z$, $t_\lambda=2\pi/(-\log\lambda)$.
The resulting flow is rotation on the circle with least period
$-\log\lambda$, which is the flow of a $\III_\lambda$ factor \cite{ConnesTakesaki}.
For $\lambda=0$ take $\Gamma=\mathbb Z+\sqrt2\mathbb Z$.
Its dual is the two-torus, and the flow is
$s\mapsto(e^{is},e^{i\sqrt2s})$.
Fourier characters show ergodicity: an invariant character
indexed by $(m,n)$ requires $m+n\sqrt2=0$, hence $m=n=0$.
Every orbit is a dense one-dimensional subgroup coset of Haar
measure zero (its bounded segments have zero area).
Thus the flow is not transitive. The Connes--Takesaki flow
classification \cite{ConnesTakesaki} gives subtype $\III_0$.
This is the same modular crossed-product mechanism as
\cite[Example 3.9]{HM}, without requiring injectivity of $N_{\rm abs}$.

For clarity non-semifiniteness is a separate conclusion.
The core flow of a semifinite factor is translation on
$L^\infty(\mathbb R)$; it is transitive and admits no invariant
normal probability state. Each of our circle or torus flows
preserves Haar probability, so cannot be that semifinite flow.
For $N_{\rm abs}$ itself non-semifiniteness was already proved by its
$S$-invariant. Hence none of the subtype conclusions silently
includes a semifinite alternative.
\end{proof}
Slofstra's commuting classification \cite{Slofstra}, also proved in \cite[Theorem 6.4]{MNY}, now proves
$\Pi^0_1$-completeness for every fixed subtype, including zero.
